\documentclass[10pt]{extarticle}
\usepackage[a4paper,margin=1.5cm]{geometry}
\usepackage[T1]{fontenc}
\usepackage{amsmath,amssymb}
\usepackage{graphicx}
\usepackage[export]{adjustbox}
\usepackage{multicol}
\usepackage{enumitem}
\usepackage{booktabs,array,tabularx}
\setlength{\columnsep}{0.7cm}
\setlength{\parindent}{0pt}
\setlength{\parskip}{3pt}
\pagestyle{empty}
\makeatletter
% figure registry (figdefs.tex is generated by the extract step)
\newcommand{\qfigdef}[4]{\@namedef{qf@file@#1}{#2}\@namedef{qf@w@#1}{#3}\@namedef{qf@wide@#1}{#4}}
\newcommand{\qf@check}[1]{\@ifundefined{qf@file@#1}{\PackageError{score2}{Unknown figure id #1}{}}{}}
\newcommand{\qf@img}[4][]{\includegraphics[width=#2\dimexpr\@nameuse{qf@w@#4}pt\relax,max width=#3,max height=0.45\textheight,#1]{\@nameuse{qf@file@#4}}}
% Full-width questions: finalize wraps every question that contains a wide figure in \qfullbegin...\qfullend,
% so the WHOLE question (stem, figures, options) is set across the page in reading order, then the two columns
% resume. The question's head - its number, stem and first figure (\qfullsplit marks the end of the head) -
% is kept on one page, so a number is never left at a page foot with its figure overleaf; the rest may break
% across pages normally, so no big gaps are left. Without a \qfullsplit, the whole question is kept together
% if it is under 45% of a page. Never write these yourself.
\newif\ifqf@full
\newif\ifqf@boxed
\newsavebox{\qf@box}
\newcommand{\qf@open}{\global\qf@fulltrue\par\global\qf@boxedtrue
  \setbox\qf@box\vbox\bgroup\hsize\textwidth\linewidth\textwidth}
\newcommand{\qf@emit}[1]{\par\egroup\global\qf@boxedfalse
  \ifdim\dimexpr\ht\qf@box+\dp\qf@box\relax<#1\textheight\noindent\usebox{\qf@box}\par\else\unvbox\qf@box\fi}
\newcommand{\qfullsplit}{\ifqf@boxed\qf@emit{0.85}\fi}
\newcommand{\qf@close}{\ifqf@boxed\qf@emit{0.45}\fi\global\qf@fullfalse}
\newcommand{\qfullbegin}{\end{multicols}\qf@open}
\newcommand{\qfullend}{\qf@close\begin{multicols}{2}}
\newcommand{\qfullstart}{\qf@open}
\newcommand{\qfulldone}{\qf@close}
% \qfig[scale]{id}: block figure, centred, directly where it is written (size set by the tool).
% Fallback if a wide figure is not inside a full-width question: break out of the columns in place - no floats.
\newcommand{\qfig}[2][1]{\qf@check{#2}%
  \ifqf@full\qf@blk{#1}{#2}\else
  \ifnum\@nameuse{qf@wide@#2}=1 \expandafter\@firstoftwo\else\expandafter\@secondoftwo\fi
  {\end{multicols}\begin{center}\qf@img{#1}{\textwidth}{#2}\end{center}\begin{multicols}{2}}%
  {\qf@blk{#1}{#2}}\fi}
\newcommand{\qf@blk}[2]{\par\smallskip{\centering\qf@img{#1}{\linewidth}{#2}\par}\smallskip}
% \qfigi[scale]{id}: inline figure (inside option lists / tables / minipages); never breaks columns.
\newcommand{\qfigi}[2][1]{\qf@check{#2}\qf@img[valign=c]{#1}{\linewidth}{#2}}
\makeatother
% structure macros - lint reads these to check completeness and figure placement
\newcommand{\papersec}[1]{\par\bigskip{\centering\Large\bfseries PAPER-#1\par}\medskip}
\newcommand{\partsec}[2]{\par\medskip\noindent{\large\bfseries PART-#1: #2}\par\smallskip}
\newcommand{\sect}[1]{\par\smallskip\noindent\textbf{#1}\par}
\newcommand{\qn}[1]{\par\smallskip\noindent\textbf{#1.}~}
\newcommand{\paraq}[2]{\par\smallskip\noindent\textbf{Paragraph for Questions #1 and #2}\par}
\newenvironment{qoptions}{\begin{itemize}[leftmargin=1.8em,itemsep=1pt,topsep=2pt,label={}]}{\end{itemize}}
\newcommand{\opt}[2]{\item[(#1)] #2}
% figure registry (generated by score2_tools.py finalize)
\qfigdef{p1_phy_s1_q01_1}{TOAS - Test 11 (27 Jun 2022)-img/p1_phy_s1_q01_1.png}{63.4}{0}
\qfigdef{p1_phy_s1_q02_1}{TOAS - Test 11 (27 Jun 2022)-img/p1_phy_s1_q02_1.png}{94.8}{0}
\qfigdef{p1_phy_s1_q03_1}{TOAS - Test 11 (27 Jun 2022)-img/p1_phy_s1_q03_1.png}{143.7}{0}
\qfigdef{p1_phy_s1_q04_1}{TOAS - Test 11 (27 Jun 2022)-img/p1_phy_s1_q04_1.png}{66.1}{0}
\qfigdef{p1_phy_s2_q06_1}{TOAS - Test 11 (27 Jun 2022)-img/p1_phy_s2_q06_1.png}{115.6}{0}
\qfigdef{p1_phy_s2_q07_1}{TOAS - Test 11 (27 Jun 2022)-img/p1_phy_s2_q07_1.png}{159.4}{0}
\qfigdef{p1_phy_s2_q08_1}{TOAS - Test 11 (27 Jun 2022)-img/p1_phy_s2_q08_1.png}{157.9}{0}
\qfigdef{p1_phy_s2_q09_1}{TOAS - Test 11 (27 Jun 2022)-img/p1_phy_s2_q09_1.png}{92.8}{0}
\qfigdef{p1_phy_s2_q10_1}{TOAS - Test 11 (27 Jun 2022)-img/p1_phy_s2_q10_1.png}{258.6}{1}
\qfigdef{p1_phy_s2_q11_1}{TOAS - Test 11 (27 Jun 2022)-img/p1_phy_s2_q11_1.png}{91.3}{0}
\qfigdef{p1_phy_s3_q02_1}{TOAS - Test 11 (27 Jun 2022)-img/p1_phy_s3_q02_1.png}{94.3}{0}
\qfigdef{p1_phy_s3_q03_1}{TOAS - Test 11 (27 Jun 2022)-img/p1_phy_s3_q03_1.png}{111.2}{0}
\qfigdef{p1_phy_s3_q06_1}{TOAS - Test 11 (27 Jun 2022)-img/p1_phy_s3_q06_1.png}{87.9}{0}
\qfigdef{p1_chem_s2_q07_1}{TOAS - Test 11 (27 Jun 2022)-img/p1_chem_s2_q07_1.png}{245.2}{0}
\qfigdef{p1_chem_s2_q08_A}{TOAS - Test 11 (27 Jun 2022)-img/p1_chem_s2_q08_A.png}{74.9}{0}
\qfigdef{p1_chem_s2_q08_B}{TOAS - Test 11 (27 Jun 2022)-img/p1_chem_s2_q08_B.png}{102.3}{0}
\qfigdef{p1_chem_s2_q08_C}{TOAS - Test 11 (27 Jun 2022)-img/p1_chem_s2_q08_C.png}{107.0}{0}
\qfigdef{p1_chem_s2_q08_D}{TOAS - Test 11 (27 Jun 2022)-img/p1_chem_s2_q08_D.png}{100.9}{0}
\qfigdef{p1_chem_s2_q09_C}{TOAS - Test 11 (27 Jun 2022)-img/p1_chem_s2_q09_C.png}{156.7}{0}
\qfigdef{p1_chem_s2_q09_D}{TOAS - Test 11 (27 Jun 2022)-img/p1_chem_s2_q09_D.png}{103.1}{0}
\qfigdef{p1_chem_s3_q02_1}{TOAS - Test 11 (27 Jun 2022)-img/p1_chem_s3_q02_1.png}{79.6}{0}
\qfigdef{p1_chem_s3_q04_1}{TOAS - Test 11 (27 Jun 2022)-img/p1_chem_s3_q04_1.png}{174.1}{0}
\qfigdef{p1_chem_s3_q05_1}{TOAS - Test 11 (27 Jun 2022)-img/p1_chem_s3_q05_1.png}{304.7}{1}
\qfigdef{p2_phy_s1_q01_1}{TOAS - Test 11 (27 Jun 2022)-img/p2_phy_s1_q01_1.png}{103.0}{0}
\qfigdef{p2_phy_s1_q02_1}{TOAS - Test 11 (27 Jun 2022)-img/p2_phy_s1_q02_1.png}{163.9}{0}
\qfigdef{p2_phy_s1_q03_1}{TOAS - Test 11 (27 Jun 2022)-img/p2_phy_s1_q03_1.png}{123.9}{0}
\qfigdef{p2_phy_s1_q04_B}{TOAS - Test 11 (27 Jun 2022)-img/p2_phy_s1_q04_B.png}{323.5}{1}
\qfigdef{p2_phy_s1_q05_1}{TOAS - Test 11 (27 Jun 2022)-img/p2_phy_s1_q05_1.png}{101.4}{0}
\qfigdef{p2_phy_s1_q06_1}{TOAS - Test 11 (27 Jun 2022)-img/p2_phy_s1_q06_1.png}{138.5}{0}
\qfigdef{p2_phy_s1_q08_1}{TOAS - Test 11 (27 Jun 2022)-img/p2_phy_s1_q08_1.png}{118.7}{0}
\qfigdef{p2_phy_s2_q09_1}{TOAS - Test 11 (27 Jun 2022)-img/p2_phy_s2_q09_1.png}{111.7}{0}
\qfigdef{p2_phy_s2_q10_1}{TOAS - Test 11 (27 Jun 2022)-img/p2_phy_s2_q10_1.png}{111.4}{0}
\qfigdef{p2_phy_s3_q01_1}{TOAS - Test 11 (27 Jun 2022)-img/p2_phy_s3_q01_1.png}{111.5}{0}
\qfigdef{p2_phy_s3_q02_1}{TOAS - Test 11 (27 Jun 2022)-img/p2_phy_s3_q02_1.png}{125.4}{0}
\qfigdef{p2_phy_s3_q03_1}{TOAS - Test 11 (27 Jun 2022)-img/p2_phy_s3_q03_1.png}{106.3}{0}
\qfigdef{p2_phy_s3_q06_1}{TOAS - Test 11 (27 Jun 2022)-img/p2_phy_s3_q06_1.png}{183.0}{0}
\qfigdef{p2_chem_s1_q03_B}{TOAS - Test 11 (27 Jun 2022)-img/p2_chem_s1_q03_B.png}{118.2}{0}
\qfigdef{p2_chem_s1_q03_C}{TOAS - Test 11 (27 Jun 2022)-img/p2_chem_s1_q03_C.png}{151.6}{0}
\qfigdef{p2_chem_s1_q03_D}{TOAS - Test 11 (27 Jun 2022)-img/p2_chem_s1_q03_D.png}{147.7}{0}
\qfigdef{p2_chem_s1_q07_1}{TOAS - Test 11 (27 Jun 2022)-img/p2_chem_s1_q07_1.png}{142.8}{0}
\qfigdef{p2_chem_s3_q02_1}{TOAS - Test 11 (27 Jun 2022)-img/p2_chem_s3_q02_1.png}{230.7}{0}
\qfigdef{p2_chem_s3_q03_1}{TOAS - Test 11 (27 Jun 2022)-img/p2_chem_s3_q03_1.png}{58.2}{0}
\qfigdef{p2_chem_s3_q04_A}{TOAS - Test 11 (27 Jun 2022)-img/p2_chem_s3_q04_A.png}{174.4}{0}
\qfigdef{p2_chem_s3_q04_B}{TOAS - Test 11 (27 Jun 2022)-img/p2_chem_s3_q04_B.png}{175.6}{0}
\qfigdef{p2_chem_s3_q05_1}{TOAS - Test 11 (27 Jun 2022)-img/p2_chem_s3_q05_1.png}{270.8}{1}
\raggedcolumns
\makeatletter\@beginparpenalty=10000\makeatother
\newcolumntype{L}[1]{>{\raggedright\arraybackslash}p{#1}}
\begin{document}
\begin{center}{\Large\bfseries TOAS --- Test 11 (27 Jun 2022)}\end{center}
\begin{multicols}{2}
\raggedcolumns
\papersec{1}
\partsec{1}{PHYSICS}
\sect{SECTION-I(i)}
\qn{1} A point charge $q$ is at a distance $r$ from the centre O of an uncharged spherical conducting layer, whose inner and outer radii are equal to $a$ and $b$ respectively. If $r<a$, then the potential at the centre O is:
\qfig{p1_phy_s1_q01_1}
\begin{qoptions}
\opt{A}{$\dfrac{q}{4\pi\varepsilon_0}\left(\dfrac1r-\dfrac1a+\dfrac1b\right)$}
\opt{B}{$\dfrac{q}{4\pi\varepsilon_0 r}$}
\opt{C}{$\dfrac{q}{4\pi\varepsilon_0}\left(\dfrac1r+\dfrac1a-\dfrac1b\right)$}
\opt{D}{$\dfrac{q}{4\pi\varepsilon_0}\left(\dfrac1a+\dfrac1b-\dfrac1r\right)$}
\end{qoptions}
\qn{2} Two wires A and C carry equal current $I_0=1$ A as shown in figure. One end of both wires extends to infinity. The magnitude of magnetic induction at point P on the angle bisector of AOC at a distance 1 cm from point O is:
\qfig{p1_phy_s1_q02_1}
\begin{qoptions}
\opt{A}{$3\times10^{-5}$ T}
\opt{B}{$6\times10^{-5}$ T}
\opt{C}{$4\times10^{-5}$ T}
\opt{D}{$8\times10^{-5}$ T}
\end{qoptions}
\qn{3} A point object O is placed in front of a sphere of radius $R$ and $\mu=\dfrac32$ at distance $\dfrac RN$. Find value of $N$ so that emergent rays are parallel to $x$-axis. Consider paraxial rays only:
\qfig{p1_phy_s1_q03_1}
\begin{qoptions}
\opt{A}{1}
\opt{B}{2}
\opt{C}{3}
\opt{D}{4}
\end{qoptions}
\qn{4} A projectile is thrown at angle $60^\circ$ with surface of earth, with velocity $v_0$ such that it attains a maximum height equal to radius of earth. Then :-
\qfig{p1_phy_s1_q04_1}
\begin{qoptions}
\opt{A}{minimum velocity would be $\dfrac{V_0}{2}$}
\opt{B}{minimum velocity would be $\dfrac{V_0}{4}$}
\opt{C}{minimum velocity would be $\dfrac{V_0}{3}$}
\opt{D}{minimum velocity would be $\dfrac{V_0}{6}$}
\end{qoptions}
\sect{SECTION-I(ii)}
\qn{5} A thin non-conducting disc of radius $R$ is uniformly charged with surface charge density $\sigma$. If potential at any point on the circumference of disc be $V$ and total electrostatic potential energy of the disc be $U$, then:
\begin{qoptions}
\opt{A}{$V=\dfrac{\sigma R}{\pi\varepsilon_0}$}
\opt{B}{$V=\dfrac{2\sigma R}{\pi\varepsilon_0}$}
\opt{C}{$U=\dfrac{2\sigma^2R^3}{3\varepsilon_0}$}
\opt{D}{$U=\dfrac43\dfrac{\sigma^2R^3}{\varepsilon_0}$}
\end{qoptions}
\qn{6} A capacitor $C$ charged to a value $Q=\varepsilon C$ is connected to a battery of emf $\varepsilon$ as shown in the figure.
\qfig{p1_phy_s2_q06_1}
\begin{qoptions}
\opt{A}{Charge on the capacitor is zero at time $RC\ln2$ after the switch S is closed.}
\opt{B}{Charge on the capacitor is zero very long time after the switch S is closed.}
\opt{C}{Maximum charge on the capacitor during the process of charging is $\varepsilon C$.}
\opt{D}{During the process of charging, the charging current first decreases and then increases exponentially.}
\end{qoptions}
\qn{7} Two capacitors of 2 $\mu$F and 3 $\mu$F are charged to 150 volt and 120 volt respectively. The plates of a capacitor are connected as shown in the figure. A discharged capacitor of capacity 1.5 $\mu$F falls to the free ends of the wire. Then after the system comes in steady state.
\qfig{p1_phy_s2_q07_1}
\begin{qoptions}
\opt{A}{Charge on the 1.5 $\mu$F capacitor will become 180 $\mu$C}
\opt{B}{Charge on the 2 $\mu$F capacitor will become 120 $\mu$C}
\opt{C}{Positive charge flows through A from left to right}
\opt{D}{Positive charge flows through A from right to left}
\end{qoptions}
\qn{8} The network shown in the figure is part of some bigger circuit. Charge on capacitor ($C=1\ \mu$F) at any time $t$ is $q=3(1-e^{-t})$ in coulombs. If at $t=0$ the potentials of point A and B is 20 V, then at $t=0$ match the following.
\qfig{p1_phy_s2_q08_1}
\begin{qoptions}
\opt{A}{Current in branch AP is 3A}
\opt{B}{Current through R is 4 A}
\opt{C}{Value of R is 4.25 $\Omega$}
\opt{D}{Value of R is 1 $\Omega$}
\end{qoptions}
\qn{9} Three very large plates are given charges as shown in the figure. If the cross-sectional area of each plate is same, then the final charge distribution on the surface of the plates a, b, c, d, e, f is:
\qfig{p1_phy_s2_q09_1}
\begin{qoptions}
\opt{A}{7Q on surface e and 3Q on surface f}
\opt{B}{$-$2Q on surface b and 3Q on surface a}
\opt{C}{$-$7Q on surface d and 2Q on surface c}
\opt{D}{The magnitude of charges at surface b, c, d, e is equal}
\end{qoptions}
\qfullbegin
\qn{10} A highly thin conducting spherical shell of radius $R_0$ is completely filled with water at $0^\circ$C. Outside surrounding temperature is $-T_0^\circ$C. Latent heat of fusion of water is $L_0$ and thermal conductivity of ice is $K_{\text{ice}}$. Assume heat is flowing symmetrically in radial direction. Choose the CORRECT option(s). Density of water is $\rho_0$. Neglect the expansion of water on freezing.
\qfig{p1_phy_s2_q10_1}
\qfullsplit
\begin{qoptions}
\opt{A}{Time in which complete water will freeze, is given by $\dfrac{\rho_0L_0R_0^2}{6K_{\text{ice}}T_0}$}
\opt{B}{Time in which complete water will freeze, is given by $\dfrac{\rho_0L_0R_0^2}{4K_{\text{ice}}T_0}$}
\opt{C}{Time in which thickness of ice will become $\dfrac{R_0}{2}$, is given by $\dfrac{\rho_0L_0R_0^2}{12K_{\text{ice}}T_0}$}
\opt{D}{Total heat released by water in freezing, is given by $\dfrac{4\pi\rho_0L_0R_0^3}{3}$}
\end{qoptions}
\qfullend
\qn{11} A parallel plate capacitor has a dielectric slab of dielectric constant $K$ between its plates that cover 1/3 of the area of its plates, as shown in the figure. The total capacitance of the capacitor is $C$ while that of the portion with dielectric in between is $C_2$. When the capacitor is charged the plane area covered by the dielectric gets charge $Q_1$ and the rest of the area gets charge $Q_2$. The electric field in the dielectric is $E_1$ and that in the other portion is $E_2$. Choose the correct option(s), ignoring edge effects.
\qfig{p1_phy_s2_q11_1}
\begin{qoptions}
\opt{A}{$\dfrac{E_1}{E_2}=K$}
\opt{B}{$\dfrac{E_1}{E_2}=\dfrac1K$}
\opt{C}{$\dfrac{Q_1}{Q_2}=\dfrac3K$}
\opt{D}{$\dfrac{C}{C_2}=\dfrac{2+K}{K}$}
\end{qoptions}
\qn{12} A particular spherical capacitor consists of two concentric conducting spherical shells of radii $R$ and $2R$. Neglect edge effects in each.
\begin{qoptions}
\opt{A}{A coaxial cylindrical capacitor, made of two concentric cylindrical conductors of radii $R$ and $2R$, has the same capacitance as the spherical capacitor. Its length should be $4R/\ln2$}
\opt{B}{A coaxial cylindrical capacitor, made of two concentric cylindrical conductors of radii $R$ and $2R$, has the same capacitance as the spherical capacitor. Its length should be $4R\ln2$}
\opt{C}{A parallel-plate capacitor, made of two circular plates of radius $R$, has the same capacitance as the spherical capacitor. The plates are at a distance of $R/8$.}
\opt{D}{A parallel-plate capacitor, made of two circular plates of radius $R$, has the same capacitance as the spherical capacitor. The plates are at a distance of $R/4$.}
\end{qoptions}
\sect{SECTION-II}
\qn{1} The current $I$ flows through an equilateral triangular loop of side length $\dfrac x8$. If magnetic field at the centroid of the triangular loop is $\dfrac{N\mu_0I}{\pi x}$, then find the value of $N/6$.
\qn{2} A body of mass 3.14 kg is suspended from one end of a wire of length 10 m. The radius of the wire is changing uniformly from $10\times10^{-4}$ m at one end to $5\times10^{-4}$ m at the other end. Find the change in length of the wire in mm. Young's modulus of the material of the wire is $2\times10^{11}$ N/m$^2$. (mass of wire is negligible) Take : $\pi=3.14$ and $g=10$ m/s$^2$
\qfig{p1_phy_s3_q02_1}
\qn{3} Calculate the reading of voltmeter (in volts) in the circuit shown below
\qfig{p1_phy_s3_q03_1}
\qn{4} A highly conducting uniform sphere of thermal capacity $C$ is heated by an electric heater, a resistance $R$ fitted within the sphere. A constant current $I$ is passed through the heater starting at time $t=0$ which gives constant power. The sphere loses heat at a rate equal to $k$ times the temperature difference between the sphere and the surrounding. The initial temperature of the sphere and that of the surrounding is $0^\circ$C. The time at which sphere attains half of its maximum attainable temperature (in $^\circ$C) is $\dfrac ck\ln a$. Then find the value of ``$a$''.
\qn{5} A body of mass 2 gm has specific heat capacity 0.25 cal/gm-$^\circ$C and surface area 30 cm$^2$. It is kept in an enclosure which is maintained at temperature 127$^\circ$C. Body is being supplied heat at rate 2 watt. If temperature of body is 27$^\circ$C and emissivity of body is 0.5. Determine the rate at which its temperature is changing (in K/s) at the given instant. (Stefan's constant is $5.67\times10^{-8}$ W/m$^2$-K$^4$. 1 cal = 4.2 J)
\qn{6} One of the circuits for the measurement of resistance by potentiometer is shown. The galvanometer is connected at point A and zero deflection is observed at length PJ = 30 cm. In second case the secondary cell is changed. Take $E_S=10$ V and $r=1\,\Omega$ in 1$^{\text{st}}$ reading and $E_S=5$ V and $r=2\,\Omega$ in 2$^{\text{nd}}$ reading. In second case, the zero deflection is observed at length PJ = 10 cm. What is the resistance $R$ (in ohm) ?
\qfig{p1_phy_s3_q06_1}
\raggedcolumns
\partsec{2}{CHEMISTRY}
\sect{SECTION-I(i)}
\qn{1} For a first order reaction A $\to$ B + C all the reactant and product molecules are in gaseous state. If the initial pressure of the reactant molecules is 100 mm and after 30 minutes the total pressure of the mixture containing molecules of all three A, B \& C gases is 140 mm, then find the time interval at which initial pressure of A will be reduced to 50 mm. ($\log_{10}3=0.5$, $\log_{10}2=0.3$)
\begin{qoptions}
\opt{A}{45 min}
\opt{B}{30 min}
\opt{C}{60 min}
\opt{D}{15 min}
\end{qoptions}
\qn{2} For the reaction:\par
2H$_2$ + 2NO $\rightleftharpoons$ N$_2$ + 2H$_2$O\par
the following mechanism has been suggested :\par
2NO $\rightleftharpoons$ N$_2$O$_2$ [Equilibrium Constant = K$_1$ (fast)]\par
N$_2$O$_2$ + H$_2$ $\xrightarrow{k_2}$ N$_2$O + H$_2$O (Slow)\par
N$_2$O + H$_2$ $\xrightarrow{k_3}$ N$_2$ + H$_2$O (fast)\par
According to above mechanism which option is correctly matched:
\begin{qoptions}
\opt{A}{$-\dfrac12\dfrac{d[\mathrm{H_2}]}{dt}=K[\mathrm{NO}][\mathrm{H_2}]$\quad Where $K=k_2\times k_3$}
\opt{B}{$-\dfrac12\dfrac{d[\mathrm{NO}]}{dt}=K[\mathrm{NO}][\mathrm{H_2}]$\quad Where $K=k_2\times K_1$}
\opt{C}{$\dfrac{d[\mathrm{N_2}]}{dt}=K[\mathrm{NO}][\mathrm{H_2}]$\quad Where $K=k_2\times K_1$}
\opt{D}{$-\dfrac12\dfrac{d[\mathrm{NO}]}{dt}=K[\mathrm{NO}]^2[\mathrm{H_2}]$\quad Where $K=k_2\times K_1$}
\end{qoptions}
\qn{3} In the cyanide extraction process of silver from argentite ore, the oxidizing and reducing agents used are
\begin{qoptions}
\opt{A}{O$_2$ and CO respectively}
\opt{B}{O$_2$ and Zn dust respectively}
\opt{C}{HNO$_3$ and Zn dust respectively}
\opt{D}{HNO$_3$ and CO respectively}
\end{qoptions}
\qn{4} Compare strength of H-bond among the following :\par
(P) H $-$ F ......... H $-$ F\par
(Q) H $-$ O ......... H $-$ O\par
(R) H $-$ N ......... H $-$ N
\begin{qoptions}
\opt{A}{P $>$ Q $>$ R}
\opt{B}{P $<$ Q $<$ R}
\opt{C}{P $>$ R $>$ Q}
\opt{D}{R $>$ P $>$ Q}
\end{qoptions}
\sect{SECTION-I(ii)}
\qn{5} Which of the following is/are true regarding FCC unit cell ? ($a$ = edge length of unit cell, $r$ = radius of lattice point)
\begin{qoptions}
\opt{A}{The smallest distance between an octahedral void and tetrahedral void is $\dfrac{\sqrt3}{4}a$.}
\opt{B}{The smallest distance between two octahedral voids is $\dfrac{a}{\sqrt2}$.}
\opt{C}{The smallest distance between two tetrahedral voids is $\dfrac{\sqrt3}{2}a$.}
\opt{D}{The smallest distance between planes of two successive layers in ABC ABC\_\_\_type packing is $\sqrt{\dfrac83}\,r$.}
\end{qoptions}
\qn{6} The wave function for an electron in H-like atom is given as $\psi(r,\theta,\phi)=\dfrac{1}{4\sqrt{2\pi}}\cdot\left(\dfrac{Z}{a_0}\right)^{5/2}\cdot r\cdot e^{-Zr/2a_0}\cdot\cos\theta$. The correct information(s) about this electron is/are -
\begin{qoptions}
\opt{A}{It belongs to a 2p orbital}
\opt{B}{This orbital has no radial node}
\opt{C}{It belongs to a 3p orbital}
\opt{D}{Total nodes for the given orbital is 1.}
\end{qoptions}
\qn{7} For the following reaction the correct option is/are :
\qfig{p1_chem_s2_q07_1}
\begin{qoptions}
\opt{A}{Both (R) \& (T) gives immediate turbidity with HCl + anhy ZnCl$_2$}
\opt{B}{(P) shows optical isomerism}
\opt{C}{(S) shows geometrical isomerism}
\opt{D}{(Q) is known as grignard reagent}
\end{qoptions}
\qn{8} Among the following pair of compounds, identify the pair(s) which can be differentiated by Fehling as well as Tollen's reagent ?
\begin{qoptions}
\opt{A}{\qfigi{p1_chem_s2_q08_A}}
\opt{B}{\qfigi{p1_chem_s2_q08_B}}
\opt{C}{\qfigi{p1_chem_s2_q08_C}}
\opt{D}{\qfigi{p1_chem_s2_q08_D}}
\end{qoptions}
\qn{9} Which of the following will give 3-pentanone.
\begin{qoptions}
\opt{A}{(CH$_3$--CH$_2$COO)$_2$Ca $\xrightarrow{\Delta}$}
\opt{B}{CH$_3$--CH$_2$--C$\equiv$N $\xrightarrow[(2)\,\mathrm{H_3O^+}]{(1)\,\mathrm{CH_3CH_2MgBr}}$}
\opt{C}{\qfigi{p1_chem_s2_q09_C}}
\opt{D}{\qfigi{p1_chem_s2_q09_D}}
\end{qoptions}
\qn{10} Select the CORRECT matches\par
\smallskip
\noindent{\footnotesize\setlength{\tabcolsep}{2pt}\begin{tabular}{@{}L{3.1cm}@{}c@{}L{4.2cm}@{}}
\multicolumn{3}{@{}c}{Column-I\hspace{1.5cm}Column-II}\\
\multicolumn{1}{@{}c}{(Change)} & & \multicolumn{1}{c}{(Given change involve)}\\[2pt]
(A) ZnO $\to$ Zn & $\Rightarrow$ & Self reduction\\[2pt]
(B) Cu(impure) $\to$ Cu(pure) & $\Rightarrow$ & Electrorefining\\[2pt]
(C) Pig iron $\to$ wrought iron & $\Rightarrow$ & Impurity removed by oxidation method\\[2pt]
(D) PbS $\to$ PbO & $\Rightarrow$ & Calcination\\
\end{tabular}}
\begin{qoptions}
\opt{A}{A}
\opt{B}{B}
\opt{C}{C}
\opt{D}{D}
\end{qoptions}
\qn{11} Which one of the following is correctly matched ?\par
\smallskip
\noindent{\footnotesize\setlength{\tabcolsep}{2pt}\begin{tabular}{@{}L{5.5cm}@{\ }L{2.9cm}@{}}
(A) 2Fe + Al$_2$O$_3\to$ 2Al + Fe$_2$O$_3$ & Thermite process for extraction of iron\\[3pt]
(B) Ti + 2I$_2\xrightarrow{523\,\mathrm{K}}$ TiI$_4\xrightarrow{1700\,\mathrm{K}}$ Ti + 2I$_2$ & Mond process\\[3pt]
(C) Ni + 4CO $\xrightarrow{\text{low temperature}}$ Ni(CO)$_4\xrightarrow{\text{high temperature}}$ Ni + 4CO & Van Arkel process\\[3pt]
(D) Cu$_2$S + 2Cu$_2$O $\to$ 6Cu + SO$_2$ & Self reduction method\\
\end{tabular}}
\begin{qoptions}
\opt{A}{(A)}
\opt{B}{(B)}
\opt{C}{(C)}
\opt{D}{(D)}
\end{qoptions}
\qn{12} Which of the following anionic radicals are identified by diluted H$_2$SO$_4$ :--
\begin{qoptions}
\opt{A}{NO$_3^-$}
\opt{B}{SO$_3^{2-}$}
\opt{C}{Cl$^-$}
\opt{D}{CO$_3^{2-}$}
\end{qoptions}
\sect{SECTION-II}
\qn{1} In a polytropic expansion process ($PV^x$ = constant), a diatomic ideal gas does an expansion work of $\dfrac y2$ Joules by absorbing `y' Joules heat from surrounding. The value of `x' is-
\qn{2} Calculate magnitude of net work in following cyclic process for 1 mole of ideal monoatomic gas (in cal.) (Given - B $\to$ C - isothermal process, D $\to$ A - adiabatic process, $T_A$ = 300 K, $T_D$ = 75 K, $R$ = 2 cal/mole-K, $\ln2=0.7$)
\qfig{p1_chem_s3_q02_1}
\qn{3} Consider all possible isomeric ketones including stereoisomers of molecular formula C$_6$H$_{12}$O, All these isomers are independently reacted with NaBH$_4$ (\textbf{NOTE} : stereoisomers are also reacted separately). The total number of ketones that give a racemic product(s) is/are
\qn{4} \qfig{p1_chem_s3_q04_1}
How many statement is/are correct for [X]\par
(a) It can give orange colour compound with 2,4--D.N.P.\par
(b) It can give positive Silver mirror test.\par
(c) It can be oxidised by Feheling solution\par
(d) It can give yellow ppt. with NaOI,$\Delta$\par
(e) It can be identified by NaHSO$_3$
\qfullbegin
\qn{5} Observe the following reaction sequence
\qfig{p1_chem_s3_q05_1}
\qfullsplit
Calculate molecular mass ($W$) of product Z and report your answer as $N$, where $N=\dfrac W3$
\qfullend
\qn{6} In Industrial extraction, the number of metals produced by reduction of their oxides with carbon or with more active metal among the following :\par
Fe, Li, Mg, Al, Au, Na, Sn, Zn
\raggedcolumns
\partsec{3}{MATHEMATICS}
\sect{SECTION-I(i)}
\qn{1} Let
\[f(x)=\begin{cases}x^2e^{2(x-1)}; & 0\le x\le1\\ a\,\mathrm{sgn}(x+1)\cos(2x-2)+bx^2, & 1<x\le2\end{cases}\]
If $f(x)$ is differentiable at $x=1$, then (where $\mathrm{sgn}(x)$ denotes signum function of $x$)
\begin{qoptions}
\opt{A}{$a=-1$, $b=2$}
\opt{B}{$a=-3$, $b=-2$}
\opt{C}{$a=-3$, $b=4$}
\opt{D}{$a=-5$, $b=-4$}
\end{qoptions}
\qn{2} The value of the definite integral $\displaystyle\int_{-\infty}^{a}\dfrac{(\sin^{-1}e^x+\sec^{-1}e^{-x})\,dx}{(\tan^{-1}e^a+\tan^{-1}e^x)(e^x+e^{-x})}$ ($a\in R$) is :
\begin{qoptions}
\opt{A}{$\ln3$}
\opt{B}{dependent on $a$}
\opt{C}{$\dfrac\pi2\ln2$}
\opt{D}{$\dfrac\pi2\ln(2\tan^{-1}e^a)$}
\end{qoptions}
\qn{3} If $y=\sin^{-1}\left(2x\sqrt{1-x^2}\right)+\sec^{-1}\dfrac{1}{\sqrt{1-x^2}}$, then $\dfrac{dy}{dx}$ can be
\begin{qoptions}
\opt{A}{$\dfrac{5}{\sqrt{1-x^2}}$}
\opt{B}{$\dfrac{2}{\sqrt{1-x^2}}$}
\opt{C}{$\dfrac{3}{\sqrt{1-x^2}}$}
\opt{D}{$\dfrac{1}{\sqrt{1+x^2}}$}
\end{qoptions}
\qn{4} From point P(0, 0) pair of tangents PA and PB are drawn to the circle $x^2+y^2+10x+10y+25=0$ and a variable tangent at a point L in minor arc AB (other than A and B) meet PA and PB at M and N respectively. The angle which MN subtends at the centre of the circle, is
\begin{qoptions}
\opt{A}{constant for all positions of L}
\opt{B}{less than $45^\circ$ for all position of L}
\opt{C}{more than $45^\circ$ for all position of L}
\opt{D}{can be less than $45^\circ$ for some positions of L}
\end{qoptions}
\sect{SECTION-I(ii)}
\qn{5} If $\displaystyle\int\sqrt{\sec^2x+3}\,dx=f(x)+\sqrt3\,g(x)+c$ where $c$ is constant of integration and $f(0)=\ln2$, $g\left(\dfrac\pi2\right)=\dfrac\pi3+1$, then
\begin{qoptions}
\opt{A}{$f(x)$ is an odd function}
\opt{B}{$f\left(\dfrac\pi4\right)=\ln\left(1+\sqrt5\right)$}
\opt{C}{$g\left(\dfrac{3\pi}{2}\right)+g\left(\dfrac\pi2\right)=2$}
\opt{D}{$g(x)$ is an odd function}
\end{qoptions}
\qn{6} If $I=\displaystyle\int_{-\pi/4}^{\pi/4}\dfrac{e^{\sec x}\,|\sin x|\sec^2x}{\left(1+e^{\sin x}\right)}\,dx$, then
\begin{qoptions}
\opt{A}{$I$ is an irrational number}
\opt{B}{$I=e^{\sqrt2}-e$}
\opt{C}{$I=3.72$}
\opt{D}{$I$ is a rational number}
\end{qoptions}
\qn{7} If $f(x)$ is differentiable at $x=5$ and $f(5)=1$, $f'(5)=\dfrac{1}{12e}$, then value of $\displaystyle\lim_{x\to5}(x-5)\int_1^{f^2(x)}\dfrac{(x+t)e^t\,dt}{(x-5)^2}$ is
\begin{qoptions}
\opt{A}{0}
\opt{B}{1}
\opt{C}{$\dfrac54$}
\opt{D}{$\dfrac56$}
\end{qoptions}
\qn{8} Let a point `R' be between P(1,0) and Q($e$,1) on the graph of $y=\ln x$ such that the tangent to the graph at point `R' is parallel to the line through points P and Q then which option is/are correct -
\begin{qoptions}
\opt{A}{$x$ co-ordinate of the point `R' is $e-1$}
\opt{B}{$x$ co-ordinate of the point `R' is $\dfrac{1}{e-1}$}
\opt{C}{if the equation of tangent at R is $y=f(x)$ then $\displaystyle\int_1^e f(x)\,dx>\int_1^e\ln x\,dx$}
\opt{D}{normal at P for $y=\ln x$ intersect $y=e^x$ at (0, 1)}
\end{qoptions}
\qn{9} If $f(x)$ be an identity function in $R$ and $g(x)=\displaystyle\sum_{k=1}^{3}\{f(x)-(2016+k)\}^{-1}$, then
\begin{qoptions}
\opt{A}{$g(x)$ is strictly decreasing in (2018, 2019)}
\opt{B}{$g(x)$ has two distinct real roots}
\opt{C}{Slope of tangent to the curve $g(x)$ at $x=f(2016)$ is $-\dfrac{49}{36}$}
\opt{D}{$\displaystyle\lim_{x\to-\infty}g(x)=0$}
\end{qoptions}
\qn{10} Let $f,g:[-2,3]\to R$ be continuous function which are twice differentiable on the interval $(-2,3)$. Let the values of $f$ and $g$ at the points $-2,0,1$ and $3$ be as given in the following table :\par
\smallskip
\noindent{\small\begin{tabular}{|c|c|c|c|c|}\hline
 & $x=-2$ & $x=0$ & $x=1$ & $x=3$\\\hline
$f(x)$ & 11 & $-4$ & 6 & 16\\\hline
$g(x)$ & 1 & $-2$ & 0 & 2\\\hline
\end{tabular}}\par
\smallskip
In each of the intervals $(-2,0)$, $(0,1)$ and $(1,3)$ the function $(f-5g)''$ never vanishes. Then the correct statement(s) is (are)
\begin{qoptions}
\opt{A}{$f'(x)-5g'(x)=0$ has exactly three solutions in $(-2,0)\cup(0,1)\cup(1,3)$}
\opt{B}{$f'(x)-5g'(x)=0$ has exactly one solutions in $(0,1)$}
\opt{C}{$f'(x)-5g'(x)=0$ has at least one solution in $(1,3)$}
\opt{D}{$f'(x)-5g'(x)=0$ has at least two solution in $(-2,0)$}
\end{qoptions}
\qn{11} Let
\[\begin{aligned}f(\theta)=\dfrac{1}{\tan^9\theta}&\left((1+\tan\theta)^{10}+(2+\tan\theta)^{10}\right.\\&\left.+\ldots+(20+\tan\theta)^{10}\right)-20\tan\theta.\end{aligned}\]
If $\displaystyle\lim_{\theta\to\frac\pi2^-}f(\theta)$ is $\ell$, then
\begin{qoptions}
\opt{A}{$\ell$ is divisible by 100}
\opt{B}{$\ell$ is divisible by 7}
\opt{C}{$\ell$ is divisible by 30}
\opt{D}{$\ell$ is a prime number}
\end{qoptions}
\qn{12} $f_1(x)=\sin^{-1}(\cos(\sin^2x))$, $f_3(x)=\sin^{-1}(\cos(\cos^2x))$, $f_2(x)=\cos^{-1}(\sin(\cos^2x))$, $f_4(x)=\cos^{-1}(\sin(\sin^2x))$ then which of the following is correct ?
\begin{qoptions}
\opt{A}{$\displaystyle\sum_{i=1}^{4}f_i(x)=2\pi-2$}
\opt{B}{$\displaystyle\sum_{i=1}^{4}(-1)^if_i(x)=0$}
\opt{C}{$\displaystyle\sum_{i=1}^{4}(-1)^{\left[\frac i3\right]}f_i(x)=0$}
\opt{D}{$f_1(x)>f_2(x)\Rightarrow\cos2x<0$}
\end{qoptions}
\sect{SECTION-II}
\qn{1} If $f:R\to R$, $f(x)=\dfrac{3x^2+mx+n}{x^2+1}$ has range $[-4,3)$ then $m^2+n^2$ is
\qn{2} The number of integral values of $k$ for which the equation $e^x=\dfrac{k}{x-3}$, has exactly two solutions is-
\qn{3} If $\displaystyle\int\dfrac{(3x^2+2x)\,dx}{x^6+2x^5+x^4+2x^3+2x^2+5}=\dfrac12\tan^{-1}(g(x))+C$, then the value of $(g(1)+g(0))$ is
\qn{4} Let $f(x)$ be a twice differentiable function satisfying $xf(x^2)+(x^2+1)f(x)=x$, $\forall\ x\in R$ then the value of $|f''(1)|$ is
\qn{5} If $f(x+1)=e^{x^2}f(x)$, where $f(x)$ is differentiable function such that $f'\left(\dfrac12\right)=5f\left(\dfrac12\right)$, then the value of $\dfrac{f'\left(N+\frac12\right)}{f\left(N+\frac12\right)}-N^2$ is ($N=1,2,\ldots\ldots$)
\qn{6} If the solution set of $\sin^{-1}x+\sin^{-1}y\ge(x+y)^2-2\left[x(y+1)-(1-y)\right]+\pi$ is $(x_0,y_0)$, then $3x_0+y_0$ is
\raggedcolumns
\newpage
\papersec{2}
\partsec{1}{PHYSICS}
\sect{SECTION-I(i)}
\qn{1} An ideal parallel plate capacitor of area $A$ is filled with three dielectric slabs having dielectric constants $K_1=3.0$, $K_2=5.0$ and $K_3=2.0$ as shown in the figure. Choose the CORRECT option(s) :-
\qfig{p2_phy_s1_q01_1}
\begin{qoptions}
\opt{A}{the equivalent capacitance of the system is $\dfrac{8\varepsilon_0A}{3d}$}
\opt{B}{the electric field in $K_1$ is more than the electric field in $K_2$.}
\opt{C}{the electric field on the left half of $K_3$ is less than the electric field on the right half of $K_3$.}
\opt{D}{the electric field on the left half of $K_3$ is more than the electric field on the right half of $K_3$.}
\end{qoptions}
\qn{2} Three capacitors of 2$\mu$F, 3$\mu$F and 5$\mu$F are individually charged with batteries of emf's 5V, 20V and 10V, respectively. After disconnecting from the voltage sources, these capacitors are connected as shown in the figure with their positive polarity plates connected to A and their negative polarity plates being earthed. Now, a battery of 20 V and an uncharged capacitor of 4 $\mu$F capacitance are connected to the junction A with a switch S, as shown. The switch is closed.
\qfig{p2_phy_s1_q02_1}
\begin{qoptions}
\opt{A}{The potential of the junction A becomes $\dfrac{100}{7}$V}
\opt{B}{Final charge on the capacitor of lowest capacitance is 200 $\mu$C}
\opt{C}{The total charge flow to the ground after closing the switch is 200 $\mu$C}
\opt{D}{The charge flow to the ground after closing the switch is zero.}
\end{qoptions}
\qn{3} A conductor having resistivity $\rho$ is bent in the shape of a half cylinder as shown in the figure. The inner and outer radii of the cylinder are $a$ and $b$ respectively and the height of the cylinder is $h$. A potential difference $V$ is applied across the two rectangular faces of the conductor using a cell of emf $V$. [Neglect the internal resistance of the cell and ignore any role of the magnetic field due to the current]
\qfig{p2_phy_s1_q03_1}
\begin{qoptions}
\opt{A}{Resistance offered by the conductor is $\dfrac{\rho\pi b}{(b-a)\,h}$.}
\opt{B}{Current flowing through the cell is $\dfrac{V\,(b-a)\,h}{\rho\pi b}$.}
\opt{C}{Current flowing through the cell is $\dfrac{V\cdot h\cdot\ln\left(\frac ba\right)}{\rho\cdot\pi}$.}
\opt{D}{Outer cylindrical surface is at a lower voltage than the inner surface.}
\end{qoptions}
\qfullbegin
\qn{4} Four long wires, with identical currents, either directly into or directly out of the page, are placed at four corners of a squares (of side $a$) in four different arrangement as shown. $B_1$, $B_2$, $B_3$ and $B_4$ are the magnitude of net magnetic field at the center of the square in respective cases.
\qfig{p2_phy_s1_q04_B}
\qfullsplit
\begin{qoptions}
\opt{A}{$B_1=B_4=0$}
\opt{B}{$B_2=\dfrac{\sqrt2\mu_0i}{\pi a}$}
\opt{C}{$B_3=\dfrac{2\mu_0i}{\pi a}$}
\opt{D}{$B_2=B_3$}
\end{qoptions}
\qfullend
\qn{5} A parallel plate air capacitor has plate area $A$ and separation between the plates $d$. Switch S is closed to connect the capacitor to a cell of emf $V$.
\qfig{p2_phy_s1_q05_1}
\begin{qoptions}
\opt{A}{Amount of heat generated in the circuit as the capacitor gets charged is $\dfrac12CV^2$}
\opt{B}{Magnitude of force that one capacitor plate exerts on the other is $\dfrac{\varepsilon_0AV^2}{d^2}$}
\opt{C}{If the distance between the plates is slowly reduced to $\dfrac d2$, the work done by external agent in the process is given by $-\dfrac12CV^2$}
\opt{D}{If the distance between the plates is reduced slowly from $d$ to $\dfrac d2$, energy dissipated in the circuit is zero.}
\end{qoptions}
\qn{6} Two rods of different materials are placed between massive walls as shown in figure. The cross section of the rods is $A$, their moduli of elasticity are $E_1$ and $E_2$ respectively. If rods are heated by $t$ degrees, then (coefficients of linear expansion of material of rods are $\alpha_1$ and $\alpha_2$ respectively)
\qfig{p2_phy_s1_q06_1}
\begin{qoptions}
\opt{A}{the strain produced in both rods is equal}
\opt{B}{the stress produced in both rods is equal}
\opt{C}{the force exerted by one rod on the other is $\dfrac{\alpha_1\ell_1+\alpha_2\ell_2}{\frac{\ell_1}{E_1}+\frac{\ell_2}{E_2}}At$}
\opt{D}{the force exerted by one rod on the other is $\dfrac{(E_1-E_2)^2A}{E_1+E_2}$}
\end{qoptions}
\qn{7} Two satellites A and B move round the earth in the same orbit. The mass of B is three times the mass of A.
\begin{qoptions}
\opt{A}{Speeds of A and B are different}
\opt{B}{The potential energy of earth $+$A is same as that of earth $+$B.}
\opt{C}{The kinetic energy of A and B are different.}
\opt{D}{The total energy of earth $+$A is same as that of earth $+$B.}
\end{qoptions}
\qn{8} On a horizontal table, there is a smooth circular groove of mean radius $R$. The walls of the groove are non conducting. Two metal balls A and B (each having mass $m$ and radius $r$) are placed inside the groove with their centers $R$ apart. The balls just fit inside the groove. The two balls are given charge $+3q$ and $-q$ are released from state of rest as shown in figure. Ignore the non-uniformity in charge distribution as the balls come close together and collide. The collision is elastic.
\qfig{p2_phy_s1_q08_1}
\begin{qoptions}
\opt{A}{Speed of ball A is just before first collision is $\sqrt{\dfrac{3q^2}{8\pi\epsilon_0m}\left(\dfrac1{2r}-\dfrac1R\right)}$}
\opt{B}{Speed of ball B is just before first collision is $\sqrt{\dfrac{3q^2}{4\pi\epsilon_0m}\left(\dfrac1{2r}-\dfrac1R\right)}$}
\opt{C}{Maximum speed acquired by each ball after they collide for the first time is $\sqrt{\dfrac{q^2}{8\pi\epsilon_0mrR}(4R-7r)}$}
\opt{D}{Maximum speed acquired by each ball after they collide for the first time is $\sqrt{\dfrac{q^2}{4\pi\epsilon_0mrR}(4R-7r)}$}
\end{qoptions}
\sect{SECTION-I(ii)}
\qn{9} \textbf{Answer the following by appropriately matching the lists based on the information given in the paragraph.}\par
There is a spherical cavity of radius $R/3$ in a spherical solid conductor of radius $R$. A is centre of solid conductor, B is centre of cavity and C is a point outside at a distance $R/3$ from the surface. In List-I, there are certain charge configuration suggested and List-II discusses the electric field and potential.
\qfig{p2_phy_s2_q09_1}
\smallskip
\noindent{\footnotesize\setlength{\tabcolsep}{2pt}\begin{tabular}{|L{.95cm}|L{2.7cm}|L{.95cm}|L{3.1cm}|}\hline
\multicolumn{2}{|c|}{List-I} & \multicolumn{2}{c|}{List-II}\\\hline
(I) & Charge $+Q$ at B & (P) & Electric field in cavity near A is to left\\\hline
(II) & Charge $+Q$ at C & (Q) & Electric field in cavity near A is to right\\\hline
(III) & Charge $-Q$ at B & (R) & Electric potential at A is $+$ve\\\hline
(IV) & Charge $-Q$ at C and $+Q$ at B & (S) & Electric potential at A is $-$ve\\\hline
 & & (T) & Electric potential at C is $+$ve\\\hline
 & & (U) & Electric potential at C is $-$ve\\\hline
\end{tabular}}
\begin{qoptions}
\opt{A}{I $\to$ P,R,T;II $\to$ R,T;III $\to$ Q,S,U;IV $\to$ P,R,U}
\opt{B}{I $\to$ P,R,T;II $\to$ R,T;III $\to$ Q,S,U;IV $\to$ Q,R,U}
\opt{C}{I $\to$ P,S,T;II $\to$ S,T;III $\to$ Q,R,U;IV $\to$ P,S,U}
\opt{D}{I $\to$ P,R,S;II $\to$ T;III $\to$ Q,S,U;IV $\to$ Q,R,U}
\end{qoptions}
\qn{10} \textbf{Answer the following by appropriately matching the lists based on the information given in the paragraph.}\par
There is a spherical cavity of radius $R/3$ in a spherical solid conductor of radius $R$. A is centre of solid conductor, B is centre of cavity and C is a point outside at a distance $R/3$ from the surface. In List-I, there are certain charge configuration suggested and List-II discusses the electric field and potential. Now a charge $+Q$ is placed at B.
\qfig{p2_phy_s2_q10_1}
\smallskip
\noindent{\footnotesize\setlength{\tabcolsep}{2pt}\begin{tabular}{|L{.95cm}|L{2.7cm}|L{.95cm}|L{3.1cm}|}\hline
\multicolumn{2}{|c|}{List-I} & \multicolumn{2}{c|}{List-II}\\\hline
(I) & Charge is moved slightly to left & (P) & Electric field at C changes.\\\hline
(II) & Charge at B is made $+2Q$ & (Q) & Electric potential at C increases.\\\hline
(III) & Charge at B is made $-Q$ & (R) & Electric potential at C decreases.\\\hline
(IV) & An external uniform electric field is switched on in the direction C to A. & (S) & Electric field at A inside cavity changes.\\\hline
 & & (T) & Electric potential at B due to charges induced on conductor changes\\\hline
 & & (U) & Electric field at B due to charges induced on conductor changes.\\\hline
\end{tabular}}
\begin{qoptions}
\opt{A}{I $\to$ P,T;II $\to$ P,Q,S,T;III $\to$ P,R,S,U;IV $\to$ P,R,S,T}
\opt{B}{I $\to$ Q,U;II $\to$ P,Q,S,U;III $\to$ P,R,T,U;IV $\to$ P,R,S,T,U}
\opt{C}{I $\to$ S,T;II $\to$ P,Q,S,T,U;III $\to$ P,R,S,T,U;IV $\to$ P,Q,T,U}
\opt{D}{I $\to$ S,U;II $\to$ P,Q,S,T;III $\to$ P,R,S,T;IV $\to$ P,R,T,U}
\end{qoptions}
\qn{11} \textbf{Answer the following by appropriately matching the lists based on the information given in the paragraph.}\par
A planet is rotating about it's axis. The radius of planet is $R$ and its mass is $m$. It has a geostationary satellite in equatorial plane at a height $h$ from surface.\par
\smallskip
\noindent{\footnotesize\setlength{\tabcolsep}{2pt}\begin{tabular}{|L{.95cm}|L{2.7cm}|L{.95cm}|L{3.1cm}|}\hline
\multicolumn{2}{|c|}{List-I} & \multicolumn{2}{c|}{List-II}\\\hline
(I) & If the planet suddenly contracts without change in it's mass & (P) & The energy of satellite decreases.\\\hline
(II) & If the planet suddenly stop rotating. & (Q) & The energy of satellite increases.\\\hline
(III) & If the geostationary satellite is pushed in direction of motion to suddenly double the velocity. & (R) & The energy of satellite remains same.\\\hline
(IV) & If the geostationary satellite suddenly detaches half of it's mass without any relative velocity. & (S) & The satellite covers larger area on the planet.\\\hline
 & & (T) & The satellite does not remain geostationary.\\\hline
 & & (U) & The satellite should have now been kept at a farther orbit from earth to be geostationary.\\\hline
\end{tabular}}
\begin{qoptions}
\opt{A}{I $\to$ R,S,T;II $\to$ R,S,T;III $\to$ Q,S,T;IV $\to$ P}
\opt{B}{I $\to$ R,S,T;II $\to$ R,S,T;III $\to$ Q,S,T;IV $\to$ Q}
\opt{C}{I $\to$ R,T,U;II $\to$ R,S,T;III $\to$ P,R,S,T;IV $\to$ R}
\opt{D}{I $\to$ R,S,T;II $\to$ R,S;III $\to$ Q,S,T;IV $\to$ P}
\end{qoptions}
\qn{12} \textbf{Answer the following by appropriately matching the lists based on the information given in the paragraph.}\par
A planet is rotating about it's axis. The radius of planet is $R$ and its mass is $m$. It has a geostationary satellite in equatorial plane at a height $h$ from surface.\par
\smallskip
\noindent{\footnotesize\setlength{\tabcolsep}{2pt}\begin{tabular}{|L{.95cm}|L{2.7cm}|L{.95cm}|L{3.1cm}|}\hline
\multicolumn{2}{|c|}{List-I} & \multicolumn{2}{c|}{List-II}\\\hline
(I) & Angular velocity of satellite. & (P) & Proportional to $\sqrt{R+h}$\\\hline
(II) & Energy of satellite & (Q) & Proportional to $\dfrac{1}{\sqrt{R+h}}$\\\hline
(III) & Angular momentum of satellite about centre of planet & (R) & Depends on mass of satellite.\\\hline
(IV) & Speed of satellite & (S) & Depends on mass of planet.\\\hline
 & & (T) & Proportional to $\sqrt G$\\\hline
 & & (U) & Depends on angular velocity of planet.\\\hline
\end{tabular}}
\begin{qoptions}
\opt{A}{I $\to$ S,T,U;II $\to$ R,S,U;III $\to$ P,R,S,T,U;IV $\to$ Q,S,T,U}
\opt{B}{I $\to$ S,T,U;II $\to$ R,S;III $\to$ P,R,S,T;IV $\to$ Q,S,T}
\opt{C}{I $\to$ R,T,U;II $\to$ R,S;III $\to$ P,R,S,T;IV $\to$ Q,T,U}
\opt{D}{I $\to$ P,S,T;II $\to$ R,S,U;III $\to$ P,R,S,T,U;IV $\to$ Q,R,S,T}
\end{qoptions}
\sect{SECTION-II}
\qn{1} An Earth satellite revolves on a circular orbit at a linear speed $v_0$. A missile device increases the absolute value of this speed to $1.5v_0$, and the satellite goes into an unknown path as shown (Fig). With what speed will the satellite leave the gravitational field of earth? The resistance of the cosmic dust is not taken into account. Mass of earth is $M$. If your answer is $\dfrac{v_0}{\alpha}$, fill $\alpha$ in OMR sheet.
\qfig{p2_phy_s3_q01_1}
\qn{2} In a meter bridge experiment, initially the jockey is at null point. Now resistance $R_1$ \& $R_2$ is interchanged. Find the value (in cm) of shift in the position of jockey.
\qfig{p2_phy_s3_q02_1}
\qn{3} Write the variation in charge ($q$) on the capacitor plates as a function of time after the switch (S) is closed at $t=0$. If $q$ is given by $q=\dfrac{CV}{2}\left[1-e^{-\frac{2t}{\alpha RC}}\right]$, then find the value of $\alpha$.
\qfig{p2_phy_s3_q03_1}
\qn{4} A spherical shell is kept in an atmosphere at temperature $T_0$. The wavelength corresponding to maximum intensity of radiation for the shell is $\lambda_0$. A point source of constant power is switched on inside the shell. The power radiated by the source is $P=0.4\,\sigma SeT_0^4$ where $S$, $e$ and $\sigma$ are outer surface area of the shell, emissivity of the outer surface of the shell and Stefan's constant respectively. Calculate the new wavelength ($\lambda$) corresponding to the maximum intensity of radiation from the shell. Assume that change in temperature ($\Delta T$) of the shell is small compared to the ambient temperature $T_0$. If value of $\lambda$ is given by $\left(\dfrac N{11}\right)\lambda_0$, then fill the value of $N$.
\qn{5} A calorimeter of water equivalent 10 gm contains a liquid of mass 50 gm at $40^\circ$C. When `$m$' gram of ice at $-10^\circ$C is put into the calorimeter and the mixture is allowed to attain equilibrium, the final temperature was found to be 20 $^\circ$C. It is known that specific heat capacity of the liquid changes with temperature as $S=\left(1+\dfrac{\theta}{500}\right)$ Cal g$^{-1}$ $^\circ$C$^{-1}$ where $\theta$ is temperature in $^\circ$C. The specific heat capacity of ice, water and colorimeter remains constant and values are $S_{\text{ice}}=0.5$ Cal g$^{-1}$ $^\circ$C$^{-1}$, $S_{\text{water}}=1.0$ Cal g$^{-1}$ $^\circ$C$^{-1}$ and latent heat of fusion of ice is $L_f=80$ Cal g$^{-1}$. Assume no heat loss to the surrounding and calculate the value of `$m$' in gram.
\qn{6} A thin equiconvex lens of glass of refractive index $\mu=\dfrac32$ and of focal length 0.3 m in air is sealed into an opening at one end of tank filled with water $\left(\mu=\dfrac43\right)$. On the opposite side of the lens, a mirror is placed inside the tank on the tank wall perpendicular to the lens axis as shown in figure. The separation between the lens and the mirror is 0.8 m. A small object is placed outside the tank in front of the lens at a distance of 0.9 m from the lens along its axis. Find the position (relative to the lens) of final image formed by the system (in cm).
\qfig{p2_phy_s3_q06_1}
\raggedcolumns
\partsec{2}{CHEMISTRY}
\sect{SECTION-I(i)}
\qn{1} For an adiabatic irreversible process involving an ideal gas against constant external pressure correct option(s) is/are :
\begin{qoptions}
\opt{A}{$\Delta U=nC_{v,m}\Delta T$}
\opt{B}{$W=nC_{v,m}\Delta T$}
\opt{C}{$W=-P_{\text{ext.}}(V_2-V_1)$}
\opt{D}{$TV^{\gamma-1}$ = constant}
\end{qoptions}
\qn{2} Which of the following reaction(s) represent standard enthalpy of formation of the product formed ?
\begin{qoptions}
\opt{A}{C(graphite) $\to$ C(diamond)}
\opt{B}{Hg ($\ell$) $\to$ Hg(g)}
\opt{C}{$\frac12$ I$_2$(solid) $+$ $\frac12$Cl$_2$(g) $\to$ ICl(s)}
\opt{D}{$\frac12$ Br$_2$(g) $+$ $\frac12$O$_2$(g) $\to$ BrO(g)}
\end{qoptions}
\qn{3} Reaction correctly matched with their major product is/are :
\begin{qoptions}
\opt{A}{Me--CO--OEt $\xrightarrow{\mathrm{NaBH_4}}$ Me--CH$_2$--OH}
\opt{B}{\qfigi{p2_chem_s1_q03_B}}
\opt{C}{\qfigi{p2_chem_s1_q03_C}}
\opt{D}{\qfigi{p2_chem_s1_q03_D}}
\end{qoptions}
\qn{4} Which of the following give 1$^\circ$ alcohol as a major product :
\begin{qoptions}
\opt{A}{H--CHO $+$ CH$_3$MgBr $\to$ $\xrightarrow{\mathrm{HOH/H^{\oplus}}}$}
\opt{B}{CH$_3$--CH$_2$--CHO $\xrightarrow{\mathrm{H_2/Pd}}$}
\opt{C}{CH$_3$--COCl $\xrightarrow[\mathrm{H_2O}]{\mathrm{NaBH_4}}$}
\opt{D}{CH$_3$--CH=CH$_2$ $\xrightarrow[\text{(ii) NaOH / H}_2\text{O}_2]{\text{(i) B}_2\text{H}_6\text{ / THF}}$}
\end{qoptions}
\qn{5} Identify correct statement (s) :
\begin{qoptions}
\opt{A}{Magnetite (Fe$_3$O$_4$) is a ferrimagnetic substance}
\opt{B}{Transition metal oxides like TiO \& CrO$_2$ show electrical conductance like metals}
\opt{C}{Solid oxide of composition Fe$_{0.95}$O is a semi-conductor}
\opt{D}{Solid MnO can behave as an antiferromagnetic substance.}
\end{qoptions}
\qn{6} In tricyclic meta phosphoric acid :
\begin{qoptions}
\opt{A}{Number of P-O-P linkage $=6$}
\opt{B}{Number of P-O-P linkage $=3$}
\opt{C}{Number of P-P linkage $=3$}
\opt{D}{Number of P-P linkage $=0$}
\end{qoptions}
\qn{7} The IUPAC name of the complex
\qfig{p2_chem_s1_q07_1}
\begin{qoptions}
\opt{A}{$\mu$-amido-$\mu$-hydroxidobis-[tetraamminecobalt(III)] ion.}
\opt{B}{Tetraammine\-cobalt(III)-\allowbreak$\mu$-amido-\allowbreak$\mu$-hydroxido\-tetraammine cobalt(III) ion}
\opt{C}{Tetraammine\-cobalt(II)-\allowbreak$\mu$-amido-\allowbreak$\mu$-hydroxido\-tetraammine cobalt(II) ion}
\opt{D}{$\mu$-amido-$\mu$-hydroxidobis-[tetraamminecobalt(II)] ion}
\end{qoptions}
\qn{8} Which of the following is/are diamagnetic and HOMO is gerade?
\begin{qoptions}
\opt{A}{N$_2$}
\opt{B}{O$_2$}
\opt{C}{O$_2^-$}
\opt{D}{C$_2$}
\end{qoptions}
\sect{SECTION-I(ii)}
\qn{9} \textbf{Answer the following by appropriately matching the lists based on the information given in the paragraph.}\par
\smallskip
\noindent{\footnotesize\setlength{\tabcolsep}{2pt}\begin{tabular}{|L{.95cm}|L{2.7cm}|L{.95cm}|L{3.1cm}|}\hline
\multicolumn{2}{|c|}{List-I} & \multicolumn{2}{c|}{List-II}\\\hline
(I) & C$_2$H$_5$Br $+$ C$_2$H$_5$Cl & (P) & $\Delta H_{\text{mix}}=0$\\\hline
(II) & CH$_3$COCH$_3$ $+$ C$_2$H$_5$OH & (Q) & $\Delta H_{\text{mix}}>0$\\\hline
(III) & CH$_3$COCH$_3$ $+$ CHCl$_3$ & (R) & $\Delta H_{\text{mix}}<0$\\\hline
(IV) & H$_2$O $+$ C$_2$H$_5$OH & (S) & can form maximum boiling azeotrope\\\hline
 & & (T) & can form minimum boiling azeotrope\\\hline
 & & (U) & $\Delta V_{\text{mix}}>0$\\\hline
\end{tabular}}\par
\smallskip
Which of the following is a \textbf{CORRECT} match ?
\begin{qoptions}
\opt{A}{IV $\to$ R, S, U}
\opt{B}{II $\to$ Q, U}
\opt{C}{III $\to$ Q, T,U}
\opt{D}{I $\to$ R, S, U}
\end{qoptions}
\qn{10} \textbf{Answer the following by appropriately matching the lists based on the information given in the paragraph.}\par
\smallskip
\noindent{\footnotesize\setlength{\tabcolsep}{2pt}\begin{tabular}{|L{.95cm}|L{2.7cm}|L{.95cm}|L{3.1cm}|}\hline
\multicolumn{2}{|c|}{List-I} & \multicolumn{2}{c|}{List-II}\\\hline
(I) & C$_2$H$_5$Br $+$ C$_2$H$_5$Cl & (P) & $\Delta H_{\text{mix}}=0$\\\hline
(II) & CH$_3$COCH$_3$ $+$ C$_2$H$_5$OH & (Q) & $\Delta H_{\text{mix}}>0$\\\hline
(III) & CH$_3$COCH$_3$ $+$ CHCl$_3$ & (R) & $\Delta H_{\text{mix}}<0$\\\hline
(IV) & H$_2$O $+$ C$_2$H$_5$OH & (S) & can form maximum boiling azeotrope\\\hline
 & & (T) & can form minimum boiling azeotrope\\\hline
 & & (U) & $\Delta V_{\text{mix}}>0$\\\hline
\end{tabular}}\par
\smallskip
Which of the following is an \textbf{INCORRECT} match ?
\begin{qoptions}
\opt{A}{II $\to$ Q, U}
\opt{B}{I $\to$ P, U}
\opt{C}{III $\to$ R, S}
\opt{D}{IV $\to$ Q, T}
\end{qoptions}
\qn{11} Match the list :\par
\smallskip
\noindent{\footnotesize\setlength{\tabcolsep}{2pt}\begin{tabular}{|L{.95cm}|L{2.7cm}|L{.95cm}|L{3.1cm}|}\hline
\multicolumn{2}{|c|}{List-I} & \multicolumn{2}{c|}{List-II}\\\hline
(I) & $[\mathrm{MnO_4}]^{-1}$ & (P) & sp$^3$\\\hline
(II) & $[\mathrm{Ni(dmg)_2}]$ & (Q) & dsp$^2$\\\hline
(III) & $[\mathrm{CrO_4}]^{2-}$ & (R) & d$^3$s\\\hline
(IV) & $[\mathrm{Zn(gly)_2}]$ & (S) & Paramagnetic\\\hline
 & & (T) & Diamagnetic\\\hline
 & & (U) & Coloured complex\\\hline
\end{tabular}}\par
\smallskip
Which of the following options has the correct combination considering List-I and List-II.
\begin{qoptions}
\opt{A}{I $\to$ Q,T,U}
\opt{B}{II $\to$ P,T,U}
\opt{C}{I $\to$ R,S,U}
\opt{D}{II $\to$ Q,T,U}
\end{qoptions}
\qn{12} Match the list :\par
\smallskip
\noindent{\footnotesize\setlength{\tabcolsep}{2pt}\begin{tabular}{|L{.95cm}|L{2.7cm}|L{.95cm}|L{3.1cm}|}\hline
\multicolumn{2}{|c|}{List-I} & \multicolumn{2}{c|}{List-II}\\\hline
(I) & $[\mathrm{MnO_4}]^{-1}$ & (P) & sp$^3$\\\hline
(II) & $[\mathrm{Ni(dmg)_2}]$ & (Q) & dsp$^2$\\\hline
(III) & $[\mathrm{CrO_4}]^{2-}$ & (R) & d$^3$s\\\hline
(IV) & $[\mathrm{Zn(gly)_2}]$ & (S) & Paramagnetic\\\hline
 & & (T) & Diamagnetic\\\hline
 & & (U) & Coloured complex\\\hline
\end{tabular}}\par
\smallskip
Which of the following options has the correct combination considering List-I and List-II.
\begin{qoptions}
\opt{A}{III $\to$ R,S,U}
\opt{B}{IV $\to$ Q,T,U}
\opt{C}{III $\to$ R,T,U}
\opt{D}{IV $\to$ P,S,U}
\end{qoptions}
\sect{SECTION-II}
\qn{1} If rate of reaction becomes double when temperature is increased from 27$^\circ$C to 37$^\circ$C. Then find activation energy of reaction in kcal/mole. ($R=2$ cal/mol-K) ($\ln2=0.7$)
\qn{2} Number of reactions in which aldehyde can be obtained as one of the product.\par
(i) CH$_3$--CH$_2$OH $\xrightarrow[300\,^\circ\mathrm{C}]{\mathrm{Cu}}$\quad (ii) CH$_3$--CH$_2$--CH$_2$--OH $\xrightarrow[\text{Pyridine}]{\mathrm{CrO_3}}$
\qfig{p2_chem_s3_q02_1}
\qn{3} Total number of moles of CO$_2$ evolved in following reaction per mole of compound `M' is -
\qfig{p2_chem_s3_q03_1}
\qn{4} For the given reactions :\par
(A) \qfigi{p2_chem_s3_q04_A}\par
(B) \qfigi{p2_chem_s3_q04_B}\par
If the molar mass of P $=a$ and the molar mass of X $=b$\par
The value of $(a-b)/2=$ ?
\qfullbegin
\qn{5} How many of the following will form stable hydrate :\par
(1) CH$_3$COCH$_3$\quad(2) CCl$_3$CHO\quad(3) CF$_3$CHO\quad(4)
\qfig{p2_chem_s3_q05_1}
\qfullend
\qn{6} Total number of stereo isomer show by complex [Rh(en)ClBrI(H$_2$O)] :
\raggedcolumns
\partsec{3}{MATHEMATICS}
\sect{SECTION-I(i)}
\qn{1} Two circles of radii $a$ and 1 touch each other externally and are inscribed in the region bounded by $y=-\sqrt{4-x^2}$ and $x$-axis, then
\begin{qoptions}
\opt{A}{$100(a+1)=125$}
\opt{B}{Number of such pairs of circles is 2}
\opt{C}{$a=1/2$}
\opt{D}{Points of contact of circle with radius $a$ and $y=-\sqrt{4-x^2}$ are $\left(\dfrac{\pm4\sqrt2}{3},-\dfrac23\right)$}
\end{qoptions}
\qn{2} Which of the following statement(s) is/are correct ?
\begin{qoptions}
\opt{A}{The point where the function changes its monotonicity and concavity is always the point of non-differentiability.}
\opt{B}{If the derivative of the function at $x=a$ is zero then $f(x)$ can be increasing or decreasing at that point.}
\opt{C}{If $f(x)$ is a twice differentiable function on R and has a relative maximum at $x=c$, then $f''(c)$ is negative.}
\opt{D}{If $f(x)$ is a differentiable function on R such that $f'(\alpha)=0$, then $f(x)$ has an extremum at $x=\alpha$ ($\alpha\in R$).}
\end{qoptions}
\qn{3} Let
\[f(x)=\begin{cases}x^3+x^2-10x; & -1\le x<0\\ \sin x; & 0\le x<\frac\pi2\\ 1+\cos x; & \frac\pi2\le x\le\pi\end{cases}\]
then $f(x)$ has
\begin{qoptions}
\opt{A}{local maximum at $x=\dfrac\pi2$}
\opt{B}{local minimum at $x=\dfrac\pi2$}
\opt{C}{absolute maximum at $x=0$}
\opt{D}{absolute maximum at $x=-1$}
\end{qoptions}
\qn{4} Let $f$ be a differentiable function on R satisfying $f'(t)=e^t(\cos^2t-\sin2t)$ and $f(0)=1$, then which of the following is/are correct ?
\begin{qoptions}
\opt{A}{$f$ is bounded in $(-\infty,0)$.}
\opt{B}{Number of solution satisfying the equation $f(t)-e^t=0$ in $[0,2\pi]$ is 3.}
\opt{C}{The value of $\displaystyle\lim_{t\to0}(f(t))^{\frac1t}=1$}
\opt{D}{$f$ is neither odd nor even.}
\end{qoptions}
\qn{5} Consider a function $f(x)=\displaystyle\int_\pi^{x^2}\dfrac{\sin x}{1+\cos^2\sqrt\theta}\,d\theta$ then which of the following is/are \textbf{CORRECT} ?
\begin{qoptions}
\opt{A}{$f(x)$ has a local minima at $x=2\pi$}
\opt{B}{$f(x)$ has a local maxima at $x=\pi$}
\opt{C}{$f(x)$ has a local minima at $x=0$}
\opt{D}{$f(x)$ has a local maxima at $x=-\pi$}
\end{qoptions}
\qn{6} Consider the function $g$ defined by
\[g(x)=\begin{cases}x^2\sin\dfrac\pi x+(x-1)^2\sin\left(\dfrac{\pi}{x-1}\right); & x\ne0,1\\ 0; & \text{if }x=0,1\end{cases}\]
then which of the following statement(s) is/are correct ?
\begin{qoptions}
\opt{A}{$g(x)$ is differentiable $\forall\ x\in R$}
\opt{B}{$g'(x)$ is discontinuous at $x=0$ but continuous at $x=1$}
\opt{C}{$g'(x)$ is discontinuous at both $x=0$ and $x=1$}
\opt{D}{Rolle's theorem is applicable for $g(x)$ in $[0,1]$}
\end{qoptions}
\qn{7} Let $f(x)=\begin{cases}(x-2)^2+3 & x\ge2\\ k+x^2 & x<2\end{cases}$, then which of the following statements is/are true ?
\begin{qoptions}
\opt{A}{If $f(x)$ is continuous at $x=2$, then $k=-1$}
\opt{B}{There exists at least one value of $k$ for which $f(x)$ is derivable $\forall\ x\in R$}
\opt{C}{If $f(x)$ has a local minima at $x=2$, then $k>-1$}
\opt{D}{If $f(x)$ is continuous at $x=2$ then $k=1$}
\end{qoptions}
\qn{8} Which of the following statement(s) is/are correct ?
\begin{qoptions}
\opt{A}{If $f$ is differentiable at $x=1$, then $\displaystyle\lim_{x\to1}\dfrac{f(x^3)-f(1)}{x-1}=3f'(1)$}
\opt{B}{If $f$ is differentiable at $x=0$, then $\displaystyle\lim_{x\to0}\dfrac{f(ax)-f(bx)}{cx}=\dfrac{a-b}{c}f'(0)$}
\opt{C}{If $f$ is differentiable at $x=x_0$, $\displaystyle\lim_{x\to x_0}\dfrac{xf(x_0)-x_0f(x)}{x-x_0}=f(x_0)-x_0f'(x_0)$}
\opt{D}{$f:R\to R$ is differentiable and $f(x)\le f(0)\ \forall\ x\in[-1,1]$ then $f'(0)=0$}
\end{qoptions}
\sect{SECTION-I(ii)}
\qn{9} \smallskip
\noindent{\footnotesize\setlength{\tabcolsep}{2pt}\begin{tabular}{|L{.9cm}|L{3.5cm}|L{.9cm}|L{1.9cm}|}\hline
\multicolumn{2}{|c|}{List-I} & \multicolumn{2}{c|}{List-II}\\\hline
(I) & $ax+by=1$ is tangent to curve $x+y=xy$, where it cuts $x$-axis, then value of $a+b$ is & (P) & $-7$\\\hline
(II) & Let $f(x)$ is continuous and differentiable over $[2,9]$ and $f'(x)\in[-4,3]$ $\forall\ x\in(2,9)$, then sum of maximum and minimum value of $f(9)-f(2)$ is & (Q) & 0\\\hline
(III) & Let $(\alpha,\beta)$ is point of inflection for the curve $y=\dfrac{x^4-40}{12}+16\ln\left(\dfrac x2\right)$, then value of $\alpha+\beta$ is & (R) & 2\\\hline
(IV) & Let $f(x)=\displaystyle\int_1^x(2t^2+3t-\lambda\log t)\,dt$ and $f'(x)$ is monotonic decreasing in interval $(0,1]$, then least value of $\lambda$ is & (S) & 7\\\hline
 & & (T) & $-2$\\\hline
 & & (U) & 1\\\hline
\end{tabular}}\par
\smallskip
Which of the following is the only correct combination ?
\begin{qoptions}
\opt{A}{I $\to$ T}
\opt{B}{I $\to$ U}
\opt{C}{II $\to$ P}
\opt{D}{II $\to$ S}
\end{qoptions}
\qn{10} \smallskip
\noindent{\footnotesize\setlength{\tabcolsep}{2pt}\begin{tabular}{|L{.9cm}|L{3.5cm}|L{.9cm}|L{1.9cm}|}\hline
\multicolumn{2}{|c|}{List-I} & \multicolumn{2}{c|}{List-II}\\\hline
(I) & $ax+by=1$ is tangent to curve $x+y=xy$, where it cuts $x$-axis, then value of $a+b$ is & (P) & $-7$\\\hline
(II) & Let $f(x)$ is continuous and differentiable over $[2,9]$ and $f'(x)\in[-4,3]$ $\forall\ x\in(2,9)$, then sum of maximum and minimum value of $f(9)-f(2)$ is & (Q) & 0\\\hline
(III) & Let $(\alpha,\beta)$ is point of inflection for the curve $y=\dfrac{x^4-40}{12}+16\ln\left(\dfrac x2\right)$, then value of $\alpha+\beta$ is & (R) & 2\\\hline
(IV) & Let $f(x)=\displaystyle\int_1^x(2t^2+3t-\lambda\log t)\,dt$ and $f'(x)$ is monotonic decreasing in interval $(0,1]$, then least value of $\lambda$ is & (S) & 7\\\hline
 & & (T) & $-2$\\\hline
 & & (U) & 1\\\hline
\end{tabular}}\par
\smallskip
Which of the following is the only correct combination ?
\begin{qoptions}
\opt{A}{III $\to$ R}
\opt{B}{III $\to$ U}
\opt{C}{IV $\to$ Q}
\opt{D}{IV $\to$ S}
\end{qoptions}
\qn{11} Let
\[f(x)=\begin{cases}x+a & ;\ x<0\\ \sqrt{x^2-2x+1} & ;\ x\ge0\end{cases}\]
and
\[g(x)=\begin{cases}x+1 & ;\ x<0\\ x^2+2(b-x)+1 & ;\ x\ge0\end{cases},\]
$a$, $b$ are non-negative real numbers. If $(g\circ f)(x)$ is continuous for all real $x$\par
\smallskip
\noindent{\footnotesize\setlength{\tabcolsep}{2pt}\begin{tabular}{|L{.9cm}|L{3.5cm}|L{.9cm}|L{1.9cm}|}\hline
\multicolumn{2}{|c|}{List-I} & \multicolumn{2}{c|}{List-II}\\\hline
(I) & The value of $b$ can be & (P) & 0\\\hline
(II) & The value of $a$ can be & (Q) & 1/2\\\hline
(III) & Value of $\dfrac{d(g(f(x)))}{dx}$ at $x=4$ is & (R) & 1\\\hline
(IV) & If $f(x)$ is non differentiable at exactly $k$ points, then the value of $k$ is & (S) & 2\\\hline
 & & (T) & 3\\\hline
 & & (U) & 4\\\hline
\end{tabular}}\par
\smallskip
Which of the following options is the only \textbf{CORRECT} combination ?
\begin{qoptions}
\opt{A}{I $\to$ R}
\opt{B}{I $\to$ Q}
\opt{C}{II $\to$ Q}
\opt{D}{II $\to$ S}
\end{qoptions}
\qn{12} Let
\[f(x)=\begin{cases}x+a & ;\ x<0\\ \sqrt{x^2-2x+1} & ;\ x\ge0\end{cases}\]
and
\[g(x)=\begin{cases}x+1 & ;\ x<0\\ x^2+2(b-x)+1 & ;\ x\ge0\end{cases},\]
$a$, $b$ are non-negative real numbers. If $(g\circ f)(x)$ is continuous for all real $x$\par
\smallskip
\noindent{\footnotesize\setlength{\tabcolsep}{2pt}\begin{tabular}{|L{.9cm}|L{3.5cm}|L{.9cm}|L{1.9cm}|}\hline
\multicolumn{2}{|c|}{List-I} & \multicolumn{2}{c|}{List-II}\\\hline
(I) & The value of $b$ can be & (P) & 0\\\hline
(II) & The value of $a$ can be & (Q) & 1/2\\\hline
(III) & Value of $\dfrac{d(g(f(x)))}{dx}$ at $x=4$ is & (R) & 1\\\hline
(IV) & If $f(x)$ is non differentiable at exactly $k$ points, then the value of $k$ is & (S) & 2\\\hline
 & & (T) & 3\\\hline
 & & (U) & 4\\\hline
\end{tabular}}\par
\smallskip
Which of the following options is the only \textbf{CORRECT} combination ?
\begin{qoptions}
\opt{A}{III $\to$ P}
\opt{B}{III $\to$ T}
\opt{C}{IV $\to$ S}
\opt{D}{IV $\to$ U}
\end{qoptions}
\sect{SECTION-II}
\qn{1} If the line $y=2-x$ is tangent to the circle $S$ at the point P(1,1) and circle $S$ is orthogonal to the circle $x^2+y^2+2x+2y-2=0$, then the length of tangent drawn from the point (2,2) to circle $S$ is
\qn{2} If the value of the integral $\displaystyle\int_{1/2}^{2}\dfrac{\cos^{-1}\left(\frac{1}{\sqrt{1+x^2}}\right)}{x}\,dx$ is equal to $\pi\ln\sqrt k$ then `$k$' is
\qn{3} Let $f$ be a differentiable function on R and satisfying $f(x)=-(x^2-x+1)e^x+\displaystyle\int_0^xe^{x-y}\cdot f'(y)\,dy$. If $f(1)+f'(1)+f''(1)=ke$, where $k\in N$, then value of $\dfrac k2$ is \underline{\hspace{1.5cm}}.
\qn{4} Let $f(x)$ be a differentiable function satisfying $f\left(\dfrac{x+y}{2}\right)=\dfrac{f(x)+f(y)}{2}$ $\forall\ x,y\in R$ and $f(0)=0$. If $\displaystyle\int_0^{2\pi}(f(x)-\sin x)^2\,dx$ is minimized, then find the value of $f(-4\pi^2)$ ?
\qn{5} If Rolle's theorem is applicable to the function $f$ defined by
\[f(x)=\begin{cases}ax^2+b & ;\ |x|<1\\ 1 & ;\ |x|=1\\ \dfrac{c}{|x|} & ;\ |x|>1\end{cases}\quad\text{for }x\in[-2,2],\]
then find the value of $(b^2+c^2+a^2)$.
\qn{6} If $\displaystyle\int(x^6+7x^5+6x^4+5x^3+4x^2+3x+1)e^x\,dx$ is equal to $\displaystyle\sum_{k=1}^{\alpha}\beta x^k\cdot e^x+C$ (where $C$ is constant of integration) then $\left(\dfrac{\alpha}{4\beta}\right)$ is
\end{multicols}
\end{document}
